% Generated draft: requires mathematical and source review.
\documentclass[11pt]{article}
\usepackage[margin=25mm]{geometry}
\usepackage{fontspec}
\setmainfont{Latin Modern Roman}
\usepackage{amsmath,amssymb,mathtools,unicode-math}
\setmathfont{Latin Modern Math}
\usepackage{enumitem,needspace,xurl,hyperref}
\setlength{\parindent}{0pt}
\setlength{\parskip}{5pt}
\setlength{\emergencystretch}{3em}
\newcommand{\N}{\mathbb{N}}
\newcommand{\Z}{\mathbb{Z}}
\newcommand{\Q}{\mathbb{Q}}
\newcommand{\R}{\mathbb{R}}
\newcommand{\C}{\mathbb{C}}
\newcommand{\F}{\mathbb{F}}
\newcommand{\A}{\mathbb{A}}
\newcommand{\PP}{\mathbb{P}}
\newcommand{\E}{\mathbb{E}}
\newcommand{\eps}{\varepsilon}
\newcommand{\dd}{\,\mathrm{d}}
\begin{document}
\section*{UnsolvedMath ID: 800790}
\textbf{Source problem code:} catalogue entry
\subsection{Definitions and mathematical statement}
Let \(n\geq 1\), and let \(\operatorname{End}(\C^n,0)\) denote the holomorphic map germs \(h:(\C^n,0)\to(\C^n,0)\) fixing the origin. A germ is superattracting if \(Dh(0)=0\). A local topological conjugacy between \(f\) and \(g\) is a homeomorphism germ \(H:(\C^n,0)\to(\C^n,0)\) satisfying \(H\circ f=g\circ H\); a local holomorphic conjugacy is such a germ that is biholomorphic.

The source imposes \(c(f)=c(g)\), but the supplied excerpt does not define \(c\). This condition is retained without assigning it an unverified meaning. Write \(\mathcal P(h)\) for the postcritical set intended by the source. For a holomorphic representative \(h\), its critical locus is
\[
\operatorname{Crit}(h)=\{z:\det Dh(z)=0\}.
\]
A usual candidate for its postcritical set is the forward union
\[
\bigcup_{j\geq 1}h^j\bigl(\operatorname{Crit}(h)\bigr),
\]
where the iterates are defined. The supplied excerpt does not specify the representative domains, whether a closure is taken, or the precise local-set convention. Consequently, this formula is explanatory, and identifying it with \(\mathcal P(h)\) requires source verification.

\textbf{Problem OP9 (Abate, 2011).} For every pair of superattracting germs \(f,g\in\operatorname{End}(\C^n,0)\) satisfying the source condition \(c(f)=c(g)\), are both of the following equivalences true?
\[
\begin{aligned}
&\exists\text{ homeomorphism germ }H:(\C^n,0)\to(\C^n,0),\quad H\circ f=g\circ H\\
&\qquad\Longleftrightarrow\quad
\exists\text{ homeomorphism germ }K:(\C^n,0)\to(\C^n,0),\quad K(\mathcal P(f))=\mathcal P(g),
\end{aligned}
\]
\[
\begin{aligned}
&\exists\text{ biholomorphism germ }H:(\C^n,0)\to(\C^n,0),\quad H\circ f=g\circ H\\
&\qquad\Longleftrightarrow\quad
\exists\text{ biholomorphism germ }K:(\C^n,0)\to(\C^n,0),\quad K(\mathcal P(f))=\mathcal P(g)?
\end{aligned}
\]
The set equalities are local at the origin. The source explicitly describes equivalence of postcritical sets by an ambient local homeomorphism or biholomorphism carrying one onto the other; it does not require this map to intertwine their dynamics. The formulation remains conditional on resolving the meanings of \(c\) and \(\mathcal P\). No finiteness or nondegeneracy hypothesis is supplied.

An appended background paragraph reports that reference [AR] contains a formal classification of superattracting germs in \(\operatorname{End}(\C^2,0)\) of order \(2\), and suggests that its methods could in principle address the general formal classification problem. It then introduces a more specific question whose statement is absent from this record. This paragraph does not supply an additional formal-conjugacy conjecture.
\subsection{Short English statement}
Abate's OP9 asks whether two superattracting holomorphic germs satisfying \(c(f)=c(g)\) are locally conjugate exactly when an ambient local map carries their postcritical sets onto one another, separately for homeomorphisms and biholomorphisms. The supplied excerpt leaves \(c\) and the precise postcritical-set convention undefined. The registry labels the problem open; that label is supplied source data, not an independently verified assessment.
\subsection{Sources}
\begin{itemize}
\item Local UnsolvedMath source registry, permanent ID \(800790\), exact source problem code \texttt{catalogue entry}, title \emph{Local Discrete Holomorphic Dynamics --- Open Problem OP9}. The record attributes the problem to Abate, with proposed year 2011 and supplied status \texttt{open}.
\item Abate, \emph{Open Problems in Local Discrete Holomorphic Dynamics} (2011), Problem OP9. Supplied source URL: \url{https://arxiv.org/abs/1107.4876}. The provenance records extraction as \texttt{arxiv\_latex} and access on 2026-10-16. This access date is supplied provenance, not a claim of independent source consultation.
\item Reference [AR], mentioned in the supplied appended paragraph as containing a formal classification of superattracting germs in \(\C^2\) of order \(2\). No bibliographic details or URL for [AR] are supplied, so its identity and the scope of that classification remain unverified.
\end{itemize}
\end{document}
